\documentclass[10pt,a4paper]{amsart}
\usepackage[margin=19mm]{geometry}
\usepackage{amsmath,amssymb,mathtools}
\usepackage[scr=rsfs,cal=euler]{mathalfa}
\usepackage{xfrac}
\usepackage[unicode,hidelinks,bookmarks=false]{hyperref}
\newcommand{\ie}{i.e.\ }
\newcommand{\cf}{cf.\ }
\newcommand{\resp}{respectively\ }
\newcommand{\Subsection}{\subsection}
\providecommand{\nobreakdash}{\nobreak\hbox{-}}
\setlength{\emergencystretch}{2em}
\multlinegap0pt
\usepackage{bm}
\usepackage{bbm}
\usepackage{dsfont}
\usepackage{xspace}
\usepackage{cancel}
\usepackage{mathtools}
\usepackage{amsthm}
\newtheorem{lemma}{Lemma}
\title[A Local Characterization of Unique Equilibrium States]{A Local Characterization of Unique Equilibrium States}
\author{C. Evans Hedges}
\date{Source: arXiv:2604.10225v2 (CC BY 4.0)}
\begin{document}
\maketitle

\noindent\emph{Source and licence.} A Local Characterization of Unique Equilibrium States. Version arXiv:2604.10225v2 (CC BY 4.0), distributed under the Creative Commons Attribution 4.0 International licence, as recorded in the arXiv licence metadata for this exact version and shown on the abstract page https://arxiv.org/abs/2604.10225v2. Full rights notice: licenses/arxiv-2604.10225-CC-BY-4.0.txt.\par\smallskip

\noindent\textit{Editorial scope.} Counted unit: the Lemma whose source label is \texttt{lem:example-continuity} in the article (source0.tex lines 450--453, source label \texttt{lem:example-continuity}), delivered together with the article's own complete original proof of it (source0.tex lines 455--488, one \texttt{proof} environment of 34 lines). No mathematics is written, repaired or re-derived: statement and proof are the article's own text line for line. Required context retained uncounted: lem:components (lines 410--420). Every definition, display and label of those lines is retained unchanged. Omitted from this package and not redistributed: 1--409, 421--449, 454--454, 489--516. Nothing inside a retained excerpt is omitted or altered: every retained body is delivered exactly as the article's lines are written, so no omission receipt of any kind accompanies this package. Citations: the retained excerpts carry no citation command at all, so no bibliography entry of the article is transcribed and no reference is redistributed. Reference adaptation: none was needed and none was made; every reference inside the delivered unit resolves inside it, so no unresolved reference or citation is produced. Printed slips kept literal: no printed word, symbol, hypothesis or number of the counted statement or of its proof was altered. Layout and class: the article's own class is \texttt{amsart} with packages including geometry, amsmath, amssymb, mathtools, needspace, etoolbox, hyperref; the wrapper is the lane's pinned \texttt{amsart} editorial wrapper at 10pt on A4, which restates the theorem-like environments the retained text needs and adds the article's own macro definitions that the retained text needs (none), with extra wrapper packages (bm, bbm, dsfont, xspace, cancel, mathtools). The delivered numbering is the unit's own. Assets: no retained span contains a figure environment, a graphics-inclusion command, a PDF-page inclusion, a table or a TikZ picture; no image file is copied into this package, the package declares no asset dependency and no image file is redistributed with it. Licence: version arXiv:2604.10225v2 of this article is distributed under the Creative Commons Attribution 4.0 International licence, as recorded in the arXiv licence metadata for this exact version and shown on the abstract page https://arxiv.org/abs/2604.10225v2. A full rights notice is delivered at licenses/arxiv-2604.10225-CC-BY-4.0.txt. Characters: every retained excerpt is pure ASCII, so the delivered engine, XeLaTeX with its default Latin Modern text font, reports no missing character.
\begin{lemma}\label{lem:components}
The topological entropy of $Y_{n,k}$ is $(\log2)/n$, and every
invariant probability $\nu$ on $Y_{n,k}$ satisfies $\nu(v)=1/n^2$.
Further, there are ergodic measures $\rho_{n,k}$ of maximal entropy on
$Y_{n,k}$ such that, for each fixed $n$,
\begin{equation}\label{eq:component-limit}
 \rho_{n,k}\longrightarrow
 \left(1-\frac{1}{n^2}\right)\delta_a+\frac{1}{n^2}\delta_b
 \qquad(k\to\infty).
\end{equation}
\end{lemma}

\begin{lemma}\label{lem:example-continuity}
The entropy map is upper semicontinuous at every ergodic measure.
It is continuous at $\delta_a$ and $\delta_b$.
\end{lemma}

\begin{proof}
We first consider an ergodic measure $\mu$ supported on
$Y=Y_{n,k}$. The second coordinate $\epsilon_{n,k}$ is isolated,
so $Y$ is both open and closed in $X$. If invariant probabilities
$\nu_j$ converge to $\mu$, then $p_j=\nu_j(Y)\to1$. For all
sufficiently large $j$, the normalized restrictions
$\eta_j=\nu_j|_Y/p_j$ are invariant probabilities on $Y$ and converge
to $\mu$. Entropy is upper semicontinuous on the finite-alphabet
subshift $Y$. Since $Y$ and its complement are invariant, affinity
and the entropy bound on $X$ give
\[
 h(\nu_j)\leq p_j h(\eta_j)+(1-p_j)\frac{\log2}{2}.
\]
Taking the upper limit gives $\limsup_j h(\nu_j)\leq h(\mu)$, as
required.

By Lemma~\ref{lem:components}, every invariant probability $\eta$
on a component $Y_{n,k}$ satisfies
\[
 h(\eta)\leq(\log2)\sqrt{\eta(v)}.
\]
This inequality also holds at the two fixed-point measures.
Ergodic decomposition and concavity of the square root therefore
give the same inequality for every invariant probability.
Since $v$ is continuous and $\delta_a(v)=0$, it follows that
$h(\nu)\to0$ as $\nu\to\delta_a$.

Finally, suppose $\nu_j\to\delta_b$. Every invariant probability
on a component $Y_{n,k}$ has $v$-average at most $1/4$, whereas
$\nu_j(v)\to1$. The total mass on these components therefore tends
to zero. Their entropies are uniformly bounded and the fixed-point
measures have entropy zero, so ergodic decomposition gives
$h(\nu_j)\to0$ as well.
\end{proof}

\end{document}
